\section{Methodology}
\label{sec:method}

\subsection{Overview of the Proposed Approach}

The method builds directly on the baseline CBMC architecture described in \Cref{sec:background}
and summarized in \Cref{fig:baseline-cbmc}. The goal is not to replace bounded model
checking, symbolic execution, or the solver stage. The goal is to construct a more accurate
checking model for OSEK/VDX applications before final verification. For ordinary C programs,
the baseline CBMC flow is often sufficient. For OSEK/VDX applications, however, the checking
model must also reflect how the deterministic scheduler, the ready queue, the task states,
and the standardized application interfaces (APIs) determine task execution order.

This section explains the proposed approach in four parts. First, it explains how the
checking model takes the OSEK/VDX scheduler into account. Second, it explains how the method
reduces unnecessary interleavings. Third, it explains how those two ideas are combined to
construct the final checking model. Fourth, it gives a short implementation note that maps
the methodology to the current realization in this branch.

Table~\ref{tab:method-comparison} provides the conceptual bridge from the baseline CBMC flow
in Section~2 to the methodology developed in the rest of this section. The comparison focuses
on the specific aspects that matter for OSEK/VDX verification: task execution order,
standardized APIs, interruption-related behaviour, unnecessary interleavings, and the
expected effect on verification accuracy.

\begin{table}[t]
  \centering
  \small
  \begin{tabularx}{\linewidth}{@{}>{\raggedright\arraybackslash}p{0.18\linewidth}>{\raggedright\arraybackslash}X>{\raggedright\arraybackslash}X@{}}
    \toprule
    Aspect & Baseline CBMC & Proposed Approach \\
    \midrule
    Task execution order & Follows ordinary program flow unless scheduler behaviour is
    encoded elsewhere. & Follows task states, the ready queue, and the deterministic
    scheduler inside the checking model. \\
    Standardized APIs & Can be treated as ordinary function effects if no OSEK/VDX-specific
    model is added. & Are interpreted as operations that may change task states, events, the
    ready queue, and therefore the next running task. \\
    Interruption-related behaviour & Broad asynchronous modelling can expose many possible
    interference points in the checking model. & Interruption-related behaviour is introduced
    only at program points where it can influence the checked execution. \\
    Unnecessary interleavings & Irrelevant interleavings may remain in the checking model and
    enlarge the verification problem. & Unnecessary interleavings are reduced before final
    bounded model checking. \\
    Expected accuracy outcome & May report a spurious bug or miss scheduler-dependent task
    execution order. & Aims to represent OSEK/VDX executions more faithfully while keeping
    CBMC's bounded model checking structure. \\
    \bottomrule
  \end{tabularx}
  \caption{Conceptual comparison between baseline CBMC and the proposed approach.}
  \label{tab:method-comparison}
\end{table}

The key point in Table~\ref{tab:method-comparison} is that the proposed approach does not
replace CBMC with a different verifier. It changes how the checking model is constructed
before final verification so that the solver is applied to a model that better reflects the
executions of OSEK/VDX applications.

To keep the discussion concrete, the remainder of this section uses one synthetic running
example. \textit{TaskControl} is a higher-priority task that waits for new sensor data.
\textit{TaskLogger} is a lower-priority task that can run while \textit{TaskControl} is
waiting. \textit{ISR\_Sensor} writes the shared sensor state and signals the event awaited by
\textit{TaskControl}. The example is intentionally small, but it is large enough to show the
three central decisions in the methodology: how task execution order is determined, how
unnecessary interleavings are reduced, and how the final checking model is built.

\subsection{Taking the OSEK/VDX Scheduler into Account}

The first methodological step is to represent the behaviour of the OSEK/VDX scheduler inside
the checking model. In an OSEK/VDX application, the next running task is not chosen by an
unstructured interleaving of threads. It is chosen by a deterministic scheduler that depends
on task states, priorities, events, and the ready queue. If the checking model ignores those
elements, the resulting task execution order no longer matches the execution order of the
application.

Listing~\ref{lst:running-example} shows the running example used throughout this section.
The example deliberately combines waiting, event signaling, shared-state update, and an
unrelated local computation so that each later methodological step can be tied back to the
same program.

\begin{lstlisting}[language=C,caption={Synthetic running example used throughout Section~3.},label={lst:running-example}]
volatile int sensor_ready = 0;
volatile int sensor_value = 0;

TASK(TaskControl)
{
  EventMaskType observed = 0;
  WaitEvent(EVENT_SENSOR_READY);
  GetEvent(TASK_ID_Control, &observed);
  if ((observed & EVENT_SENSOR_READY) != 0u && sensor_ready) {
    int sample = sensor_value;
    assert(sample >= 0);
  }
  ClearEvent(EVENT_SENSOR_READY);
  TerminateTask();
}

TASK(TaskLogger)
{
  int checksum = 0;
  checksum = checksum + 1;
  checksum = checksum * 2;
  TerminateTask();
}

ISR(ISR_Sensor)
{
  sensor_value = 7;
  sensor_ready = 1;
  SetEvent(TASK_ID_Control, EVENT_SENSOR_READY);
}
\end{lstlisting}

Table~\ref{tab:running-config} summarizes the intended execution context for this example.
The details are simple on purpose: \textit{TaskControl} begins as the running task,
\textit{TaskLogger} is already ready at a lower priority, and the interruption can make the
higher-priority task ready again by signaling the awaited event.

\begin{table}[t]
  \centering
  \small
  \begin{tabularx}{\linewidth}{@{}>{\raggedright\arraybackslash}p{0.22\linewidth}>{\raggedright\arraybackslash}p{0.22\linewidth}X@{}}
    \toprule
    Entity & Priority / type & Initial status and role \\
    \midrule
    \textit{TaskControl} & Higher-priority extended task & Running at the beginning of the
    path; waits for \texttt{EVENT\_SENSOR\_READY}; later reads the event state and the shared
    sensor value. \\
    \textit{TaskLogger} & Lower-priority task & Ready while \textit{TaskControl} is running;
    may execute while \textit{TaskControl} is waiting; performs only local computation in
    this example. \\
    \textit{ISR\_Sensor} & Interruption & Writes \texttt{sensor\_ready} and
    \texttt{sensor\_value}, then signals the event awaited by \textit{TaskControl}. \\
    \bottomrule
  \end{tabularx}
  \caption{Configuration summary for the running example.}
  \label{tab:running-config}
\end{table}

The central point is that the standardized APIs in
Listing~\ref{lst:running-example} cannot be treated as ordinary function effects. When
\textit{TaskControl} executes \texttt{WaitEvent}, the important effect is not simply that a
call occurs. The important effect is that the task may move from the running state to the
waiting state. Once that happens, \textit{TaskLogger} may become the next runnable task.
Later, when \textit{ISR\_Sensor} executes \texttt{SetEvent}, the event state changes and
\textit{TaskControl} may become ready again. Because \textit{TaskControl} has higher
priority, the deterministic scheduler may then dispatch it ahead of \textit{TaskLogger}.
The checking model must therefore represent state changes in the scheduler domain, not merely
ordinary control-flow changes in the current C function.

Table~\ref{tab:running-trace} is the output of this first methodological step. It shows how
the same source program leads to different verification states once the scheduler, the ready
queue, and the task states are represented explicitly in the checking model.

\begin{table}[t]
  \centering
  \footnotesize
  \begin{tabularx}{\linewidth}{@{}>{\raggedright\arraybackslash}p{0.07\linewidth}>{\raggedright\arraybackslash}p{0.20\linewidth}>{\raggedright\arraybackslash}p{0.23\linewidth}>{\raggedright\arraybackslash}p{0.16\linewidth}>{\raggedright\arraybackslash}p{0.11\linewidth}>{\raggedright\arraybackslash}X@{}}
    \toprule
    Step & Current context & Task states & Ready queue & Event state & Next dispatch decision \\
    \midrule
    0 & Start of path & \textit{TaskControl}: running; \textit{TaskLogger}: ready & \textit{TaskLogger} & Clear & Keep \textit{TaskControl} running \\
    1 & \textit{TaskControl} executes \texttt{WaitEvent} while the event is clear & \textit{TaskControl}: waiting; \textit{TaskLogger}: ready & \textit{TaskLogger} & Clear & Switch to \textit{TaskLogger} \\
    2 & \textit{TaskLogger} executes local computation & \textit{TaskControl}: waiting; \textit{TaskLogger}: running & Empty & Clear & Keep \textit{TaskLogger} running unless interrupted \\
    3 & \textit{ISR\_Sensor} writes shared state and signals the event & \textit{TaskControl}: ready; \textit{TaskLogger}: ready or preemptable & \textit{TaskControl}, \textit{TaskLogger} & Set & Preempt with \textit{TaskControl} \\
    4 & \textit{TaskControl} resumes, observes the event, and reads shared data & \textit{TaskControl}: running; \textit{TaskLogger}: ready & \textit{TaskLogger} & Set, then cleared & Keep \textit{TaskControl} running \\
    \bottomrule
  \end{tabularx}
  \caption{State evolution for the running example when the checking model tracks task states,
  the ready queue, and the event state explicitly.}
  \label{tab:running-trace}
\end{table}

Figure~\ref{fig:osek-sequence} summarizes the same causal chain at a higher level. The
figure is not a control-flow graph. It is a compact view of how a standardized API changes
task states and events, how those changes reach the ready queue, and how the deterministic
scheduler turns that information into the next task execution step in the checking model.

\begin{figure}[t]
  \centering
  \begin{tikzpicture}[
  >=Latex,
  font=\small,
  box/.style={
    draw,
    rounded corners=2pt,
    align=center,
    text width=24mm,
    minimum height=9mm,
    inner sep=3pt,
    fill=orange!6
  },
  arrow/.style={->, thick}
]
  \node[box] (api) at (0,2.05) {Standardized\\APIs};
  \node[box] (state) at (3.25,1.05) {Task states\\and events};
  \node[box] (queue) at (3.25,-1.05) {Ready\\queue};
  \node[box] (sched) at (0,-2.05) {Deterministic\\scheduler};
  \node[box] (order) at (-3.25,-1.05) {Task execution\\order};
  \node[box] (step) at (-3.25,1.05) {Next checking\\step};

  \draw[arrow] (api) -- (state);
  \draw[arrow] (state) -- (queue);
  \draw[arrow] (queue) -- (sched);
  \draw[arrow] (sched) -- (order);
  \draw[arrow] (order) -- (step);
  \draw[arrow] (step) -- (api);
\end{tikzpicture}

  \caption{Conceptual cycle from standardized application interfaces (APIs) to task states,
  the ready queue, the deterministic scheduler, and the next checking step.}
  \label{fig:osek-sequence}
\end{figure}

Read clockwise, the cycle says the following. A standardized API such as \texttt{WaitEvent},
\texttt{GetEvent}, or \texttt{SetEvent} changes the scheduler-visible state of the
application. Those changes update the task states and the event information associated with
the affected tasks. The resulting state determines the contents of the ready queue. The ready
queue is then interpreted by the deterministic scheduler to determine task execution order.
That task execution order determines the next checking step, which may in turn lead to
another standardized API or another scheduler-relevant state update. Baseline CBMC follows
ordinary control-flow effects, whereas the proposed approach must also update task state and
scheduler state in the checking model.

In the running example, this means that the verification path should not simply continue with
the syntactically next statement after \texttt{WaitEvent}. Instead, the checking model should
represent that \textit{TaskControl} is blocked, that \textit{TaskLogger} may run while the
higher-priority task is waiting, and that \textit{ISR\_Sensor} can make
\textit{TaskControl} ready again. Only then does the task execution order in the
verification result reflect the execution order defined by the OSEK/VDX scheduler.

\subsection{Reducing Unnecessary Interleavings}

Taking the scheduler into account is not sufficient by itself. The second methodological step
is to reduce unnecessary interleavings in the checking model. Existing checking approaches for
concurrent software often consider a very broad set of asynchronous behaviours. In the
present setting, that can enlarge the verification problem and can also produce a spurious
bug when interruption-related behaviour is introduced at program points that are irrelevant to
the property being checked.

The key question is therefore not ``where could an interruption happen in general?'' The key
question is ``where can interruption-written state affect the checked execution?'' This is a
narrower and more useful question because it focuses the checking model on observable effects.
If interruption-written state cannot influence a later branch, a task-state transition, or a
property-relevant shared-data read, then representing interruption-related behaviour at that
point adds interleavings without improving accuracy.

The running example makes this distinction concrete. \textit{ISR\_Sensor} writes the event
state observed by \textit{TaskControl} and the shared variables \texttt{sensor\_ready} and
\texttt{sensor\_value}. Those writes matter near the points where \textit{TaskControl}
decides whether it should block, where it observes the signaled event, where it branches on
that information, and where it reads the shared sensor value. By contrast, the two checksum
updates inside \textit{TaskLogger} are purely local in this example. They do not observe the
event state, they do not change the ready queue, and they do not influence the later branch
or shared-data read that matter for the checked path.

Table~\ref{tab:selected-points} shows the output of this second methodological step. The
table is intentionally human-readable. It does not list raw artifact fields. It records only
the conceptual question answered by the analysis: which program points must retain
interruption-related behaviour because they can affect the checked execution?

\begin{table}[t]
  \centering
  \footnotesize
  \begin{tabularx}{\linewidth}{@{}>{\raggedright\arraybackslash}p{0.33\linewidth}>{\raggedright\arraybackslash}p{0.19\linewidth}>{\raggedright\arraybackslash}p{0.12\linewidth}X@{}}
    \toprule
    Program point in the running example & Depends on interruption-written state? & Keep / discard & Reason \\
    \midrule
    Immediately before \texttt{WaitEvent(EVENT\_SENSOR\_READY)} & Yes & Keep & The interruption can determine whether the task remains running or moves to the waiting state. \\
    Immediately before \texttt{GetEvent(TASK\_ID\_Control, \&observed)} & Yes & Keep & The returned event mask depends on whether the interruption has already signaled the awaited event. \\
    Immediately before the branch that checks \texttt{observed} and \texttt{sensor\_ready} & Yes & Keep & The branch condition depends directly on state written by the interruption. \\
    Immediately before the read of \texttt{sensor\_value} & Yes & Keep & The shared value read by \textit{TaskControl} is produced by the interruption. \\
    Immediately before \texttt{checksum = checksum + 1} in \textit{TaskLogger} & No & Discard & The computation is local to \textit{TaskLogger} and cannot observe interruption-written state in this example. \\
    Immediately before \texttt{checksum = checksum * 2} in \textit{TaskLogger} & No & Discard & This step is also local and does not affect task release, event state, or the later shared-data read. \\
    \bottomrule
  \end{tabularx}
  \caption{Conceptual selection result for interruption-relevant program points in the
  running example.}
  \label{tab:selected-points}
\end{table}

The point of Table~\ref{tab:selected-points} is not to predict the exact moment of every
interruption. The point is to keep only those program points where interruption-written state
can change the path explored by bounded model checking. Baseline CBMC with broad
asynchronous modelling may still explore both the kept and the discarded points. The
proposed approach keeps the first group and removes the second group from the final checking
model, which reduces unnecessary interleavings without discarding the scheduler-relevant
behaviour.

\subsection{Building the Final Checking Model}

The proposed approach combines the previous two ideas by separating two decisions that are
often mixed together. The first decision is analytical: identify the program points at which
interruption-related behaviour can influence the checked execution. The second decision is
constructive: build a new checking model that introduces interruption-related behaviour only
at those selected points.

This separation matters because the two decisions answer different questions. The analytical
step asks where interruption-written state can affect task execution order or a
property-relevant computation. The constructive step asks how to represent those effects in
the final checking model. Keeping those steps separate avoids a monolithic transformation in
which relevance and model rebuilding are entangled. As a result, the final checking model can
remain narrower than a blanket insertion of interruption-related behaviour at many unrelated
points.

For the running example, the analytical result can be summarized in the compact form shown in
Listing~\ref{lst:selected-summary}. This summary is the bridge between the two stages: it is
already specific enough to guide reconstruction, but it is still expressed in terms of the
checked path rather than in low-level artifact details.

\begin{lstlisting}[caption={Human-readable summary of the selected program points for the running example.},label={lst:selected-summary}]
Selected points
1. Before WaitEvent(EVENT_SENSOR_READY)
2. Before GetEvent(TASK_ID_Control, &observed)
3. Before the branch that checks observed and sensor_ready
4. Before the read of sensor_value

Excluded points
- The local checksum updates inside TaskLogger
\end{lstlisting}

Figure~\ref{fig:improved-pipeline} shows how the final checking model is then built. The
upper lane shows the baseline CBMC path. The lower lane shows the additional construction
stages used by the proposed approach before final bounded model checking.

\begin{figure}[t]
  \centering
  \begin{tikzpicture}[
  node distance=7mm and 6mm,
  box/.style={
    draw,
    rounded corners=2pt,
    align=center,
    text width=30mm,
    minimum height=9mm,
    fill=blue!4
  },
  arrow/.style={-{Latex[length=2mm]},thick}
]
  \node[box] (src) {Source program\\and configuration\\data};
  \node[box,right=of src] (gotocc) {Initial checking\\model};
  \node[box,right=of gotocc] (analysis) {Interleaving\\relevance\\analysis};
  \node[box,below=of analysis] (manifest) {Selected\\program points};
  \node[box,left=of manifest] (aib) {Restricted-model\\construction};
  \node[box,left=of aib] (cbmc) {Final checking\\model for bounded\\model checking};

  \draw[arrow] (src) -- (gotocc);
  \draw[arrow] (gotocc) -- (analysis);
  \draw[arrow] (analysis) -- (manifest);
  \draw[arrow] (manifest) -- (aib);
  \draw[arrow] (aib) -- (cbmc);
\end{tikzpicture}

  \caption{Baseline and improved checking-model construction. Numbered blocks indicate stages
  added by the proposed approach.}
  \label{fig:improved-pipeline}
\end{figure}

Conceptually, the upper lane starts from the original program and proceeds directly to model
construction, symbolic execution, and bounded model checking. The lower lane inserts four
additional responsibilities. The first added block builds a project-scoped view of the
program rather than a single straight-line execution view. The second added block identifies
where interruption-written state can affect the checked path. The third added block rebuilds
a restricted source-level checking model that retains interruption-related behaviour only at
those selected points. The fourth added block applies scheduler-consistent symbolic execution
so that standardized APIs change task states, the ready queue, and task execution order in a
way that matches OSEK/VDX semantics. This is the second major difference from baseline CBMC:
the proposed approach changes both where interruption-related behaviour appears and how the
resulting task execution order is interpreted during verification.

Listing~\ref{lst:restricted-checking-model} shows the output of the constructive step for the
running example. The guarded calls are part of the checking model, not part of the original
application program. They are placed only at the selected points. No such call is inserted
inside the local checksum updates of \textit{TaskLogger}, because those updates were excluded
as unnecessary interleavings in Table~\ref{tab:selected-points}.

\begin{lstlisting}[language=C,caption={Illustrative restricted checking model for the running example.},label={lst:restricted-checking-model}]
TASK(TaskControl)
{
  EventMaskType observed = 0;
  if (nondet_bool()) ISR_Sensor(0);
  WaitEvent(EVENT_SENSOR_READY);
  if (nondet_bool()) ISR_Sensor(0);
  GetEvent(TASK_ID_Control, &observed);
  if (nondet_bool()) ISR_Sensor(0);
  if ((observed & EVENT_SENSOR_READY) != 0u && sensor_ready) {
    if (nondet_bool()) ISR_Sensor(0);
    int sample = sensor_value;
    assert(sample >= 0);
  }
  ClearEvent(EVENT_SENSOR_READY);
  TerminateTask();
}

TASK(TaskLogger)
{
  int checksum = 0;
  checksum = checksum + 1;
  checksum = checksum * 2;
  TerminateTask();
}
\end{lstlisting}

The methodological role of Listing~\ref{lst:restricted-checking-model} is straightforward.
The checking model still admits interruption-related behaviour, but only at points where that
behaviour can change the execution being checked. This preserves the main advantage claimed
in Sections~1 and~2: fewer unnecessary interleavings without ignoring the OSEK/VDX
scheduler or the shared state written by interruptions.

\subsection{Implementation Note}

The current branch realizes this methodology with a small set of focused extensions rather
than a monolithic verifier rewrite. Table~\ref{tab:implementation-note} summarizes the role
of each concrete component and the artifact or checking-model effect it produces.

\begin{table}[t]
  \centering
  \small
  \begin{tabularx}{\linewidth}{@{}>{\raggedright\arraybackslash}p{0.26\linewidth}>{\raggedright\arraybackslash}X>{\raggedright\arraybackslash}X@{}}
    \toprule
    Component & What it consumes and does & What it produces \\
    \midrule
    \shortstack[l]{\texttt{goto-cc}\\extension} & Consumes the source program and configuration data, then
    builds the initial checking model while preserving the project-level view needed by the
    later analysis stage. & An initial checking model together with the project-scoped
    information needed to analyze interruption relevance. \\
    \shortstack[l]{\texttt{interleaving-}\\\texttt{analysis}} & Consumes the initial checking model and the designated
    interruption sources, then identifies program points where interruption-written state can
    affect the checked path. & A selection artifact describing the retained program points for
    restricted reconstruction. \\
    \shortstack[l]{\texttt{aib} with\\\texttt{lib-interleaving-block}} & Consumes the original project tree
    and the selection artifact, then rebuilds a restricted source-level checking model with
    guarded interruption calls only at the retained points. & A rewritten project view that is
    narrower than blanket interruption insertion. \\
    \shortstack[l]{\texttt{goto-symex} with\\\texttt{os-api}} & Consumes the restricted checking model and the
    OSEK/VDX configuration, then realizes the effects of standardized APIs on task states,
    events, the ready queue, and dispatch decisions during symbolic execution. & Scheduler-
    consistent symbolic states for final bounded model checking. \\
    \bottomrule
  \end{tabularx}
  \caption{Concrete realization of the methodology in the current branch.}
  \label{tab:implementation-note}
\end{table}

Figure~\ref{fig:improved-pipeline} uses these component boundaries implicitly. The added
blocks in the lower lane correspond, in order, to the project-scoped model-construction
extension around \texttt{goto-cc}, the separate \texttt{interleaving-analysis} stage, the
\texttt{aib} rewriting step built on \texttt{lib-interleaving-block}, and the
\texttt{goto-symex}/\texttt{os-api} realization of OSEK/VDX API effects during symbolic
execution.

Exact command-line spellings, manifest-field names, and concrete artifact formats are kept
out of the methodology body and collected in \Cref{sec:appendix}. That appendix is the right
place for interface-level details once the theoretical flow in this section is understood.
